\section{Related Work}
\label{sec:related}

\subsection{Multi-Task Model Merging}
Model merging has emerged as a key paradigm for combining multiple specialized networks into a single multi-task system without incurring the costs of joint training or multi-task fine-tuning \cite{choshen2022fusing}. A seminal approach is \emph{Task Arithmetic} \cite{ilharco2022editing}, which constructs "task vectors" by subtracting the pre-trained base model weights from the fine-tuned expert weights, then linearly combines these vectors and adds them back to the base model. This simple additive approach has shown remarkable efficacy in combining diverse modalities and tasks \cite{wortsman2022robust}. 
However, uniform scaling of task vectors assumes that all layers contribute equally to each task. To relax this assumption, several static, non-uniform merging methods have been proposed. For example, Fisher Merging \cite{matena2022merging} scales parameter updates proportional to their diagonal Fisher Information. RegCalMerge \cite{regcalmerge2026} leverages class-capacity normalization and elastic spatial regularization to mitigate sacrificial task bias. While static non-uniform methods improve merging accuracy, their coefficients remain fixed post-merging, rendering them unable to adapt dynamically to covariate shifts or localized test-time distribution changes.

\subsection{Test-Time Adaptation and Dynamic Merging}
Test-Time Adaptation (TTA) aims to adjust a pre-trained model to a shifting target domain on-the-fly, using only unlabeled streaming test data \cite{wang2020tent, liang2023comprehensive}. A prominent objective in TTA is the self-supervised minimization of Shannon entropy \cite{wang2020tent}, which shifts decision boundaries away from high-density test regions. Recently, \emph{AdaMerging} \cite{lu2023adamerging} pioneered the integration of TTA with model merging. Instead of updating the massive weight matrices of the model directly, AdaMerging keeps the expert weights frozen and adaptively optimizes layer-wise merging coefficients ($\lambda_{k, l}$) on the unlabeled test stream.
Despite its potential, we expose that unconstrained layer-wise optimization at test-time suffers from the \emph{Overfitting-Optimizer Paradox}. Because TTA relies on small local batches (often of size $16$ to $64$), the empirical entropy loss is heavily corrupted by sampling noise. An unconstrained optimizer with $K \times L$ free variables quickly overfits to these local transductive fluctuations, memorizing high-frequency batch noise and causing the coefficients to deviate wildly from the true global sensitivity profile of the network. This results in \emph{representation collapse}---where the model's general capability severely degrades. 

\subsection{Continuous Parameterization and Numerical Subspaces}
To combat transductive overfitting, researchers have sought to restrict the optimization degrees of freedom. A notable attempt is \emph{PolyMerge} \cite{polymerge2024}, which represents layer-wise coefficients as a continuous function of network depth, parameterized via a low-degree polynomial. By limiting the polynomial degree $d \ll L$, PolyMerge acts as an implicit low-pass filter that rejects high-frequency transductive noise, significantly improving generalization.
However, PolyMerge parameterizes these continuous curves using the standard power basis (monomials $\bar{l}^j$), which are notorious in approximation theory for their severe numerical ill-conditioning \cite{gautschi1990influence}. The resulting design matrix is Vandermonde-type, and the condition number of its corresponding Gram matrix grows exponentially with the polynomial degree $d$. This ill-conditioning distorts the optimization landscape, creating extremely steep, asymmetric valleys that slow down gradient-based optimizers, induce vanishing gradients for higher-order terms, and compromise convergence.

In contrast, \textbf{ChebyMerge} utilizes an orthogonal basis spanned by Chebyshev polynomials of the first kind ($T_j(x)$). By leveraging discrete orthogonality, we maintain near-perfect numerical conditioning regardless of the degree, enabling rapid and isotropic convergence while preserving all the noise-filtering and generalization benefits of a continuous low-degree subspace.
